% house_prelim.tex: HOUSE COPY of definite_integrals/common/prelim_style.tex; [HOUSE] titles from \housechapter
% =====================================================================
% prelim_style.tex: makes the Test look like a Singapore JC H2 Mathematics
% 9758 prelim paper (studied from ACJC/RI 2026 P1/P2 on the box:
% body Times New Roman + Times-style maths, cover in Arial).
% Fonts: prelim_fonts.tex (real Times New Roman/Arial + Times-style maths).
% Shared by Ch10 (di_style) and Ch11 (de_style; sets \prelimtitle/\prelimsubject). Loaded AFTER the chapter style (overrides its layout for the test only).
% =====================================================================
% [HOUSE] fonts loaded by house.tex
\providecommand{\prelimtitle}{\MakeUppercase{\housechap} TEST}
\providecommand{\prelimsubject}{\housechap}
\geometry{a4paper, top=22mm, bottom=24mm, left=25mm, right=25mm, headsep=6mm, footskip=12mm}
\setlength{\parindent}{0pt}
\setlength{\parskip}{0.9ex plus 0.3ex minus 0.2ex}
\linespread{1.3}
\setlength{\emergencystretch}{3em}
\clubpenalty=10000 \widowpenalty=10000 \interlinepenalty=10000
\AtBeginDocument{\fontsize{12}{14.4}\selectfont}
\renewcommand{\normalsize}{\fontsize{12}{14.4}\selectfont}

% ---- figures: plain black/grey like a printed paper ----------------------
\ifdefined\prelimkeepcolour\else % question paper: black/grey figures; solutions keep colour
\pgfplotsset{curve/.style={black, semithick, smooth, samples=150},
  line2/.style={black, semithick, samples=2},
  sketch/.append style={axis line style={thin, black}, every axis label/.append style={black}}}
\tikzset{region/.style={fill=black!15, draw=none}, regionB/.style={fill=black!30, draw=none},
  rect/.style={fill=black!8, draw=black, thin}, ax3/.style={-{Stealth[length=1.5mm]}, thin, black}}
\colorlet{blue}{black}\colorlet{red}{black}\colorlet{orange}{black}
\fi

% ---- header / footer: page number centred top (bold), [Turn over at bottom right
\makeatletter
\providecommand{\prelimfootlabel}{\prelimtitle\ 2026}
\newcommand{\pp@footer}{{\coverfont\scriptsize \prelimfootlabel\hspace{3em}H2 MATHEMATICS 9758}}
\fancypagestyle{prelim}{\fancyhf{}%
  \chead{\bfseries\thepage}%
  \lfoot{\pp@footer}%
  \rfoot{\ifnum\value{page}<\getpagerefnumber{LastPage}\relax\textbf{[Turn over}\fi}%
  \renewcommand{\headrulewidth}{0pt}\renewcommand{\footrulewidth}{0pt}}
\fancypagestyle{prelimcover}{\fancyhf{}%
  \rfoot{\textbf{[Turn over}}%
  \renewcommand{\headrulewidth}{0pt}\renewcommand{\footrulewidth}{0pt}}
\pagestyle{prelim}

% ---- marks: [n] regular weight, right-aligned at the end of the part ----------
\renewcommand{\mk}[1]{{\unskip\nobreak\hfil\penalty50\hskip1.5em\hbox{}\nobreak\hfill[#1]%
  \parfillskip=0pt\finalhyphendemerits=0\par}}

% ---- questions: bold number at the left margin, text indented; grey [tag] straight after the number (3B)
\newsavebox{\pp@qmeas}
\newcommand{\pp@qlist}[1]{%
  \begin{list}{}{\setlength\leftmargin{2.6em}\setlength\labelwidth{2.1em}\setlength\labelsep{0.5em}%
     \setlength\topsep{0pt}\setlength\itemsep{0pt}\setlength\parsep{0.9ex}%
     \renewcommand\makelabel[1]{\textbf{##1}\hfil}}%
  \item[\thediq]\@endparpenalty\@M\setbox0=\hbox{\di@tag{#1}}\ifdim\wd0=0pt\else{\small\textcolor{black!60}{[\di@tag{#1}]}}\par\nobreak\fi\csname diQbody@#1\endcsname\par\penalty\@M
  \end{list}}
% a question starts on a new page unless it fits whole on the current one
% (questions taller than a page may break, but only between parts)
\renewcommand{\putQ}[1]{\stepcounter{diq}%
  \expandafter\xdef\csname diLab@#1\endcsname{\thediq}%
  \par\setbox\pp@qmeas=\vbox{\pp@qlist{#1}}%
  \ifdim\dimexpr\ht\pp@qmeas+\dp\pp@qmeas\relax<0.95\textheight
    \needspace{\dimexpr\ht\pp@qmeas+\dp\pp@qmeas\relax}\else\needspace{14\baselineskip}\fi
  \zsavepos{q-s\thediq}\pp@qlist{#1}\par\zsavepos{q-e\thediq}\vspace{5ex plus 2ex minus 1ex}}
% part labels (a)/(i) bold, indented under the question text
\renewenvironment{qparts}[1][2.6em]{\begin{list}{}{\setlength\leftmargin{#1}%
  \setlength\labelwidth{\dimexpr#1-0.5em\relax}\setlength\labelsep{0.5em}%
  \setlength\itemsep{1.4ex}\setlength\topsep{1.2ex}\setlength\parsep{0.8ex}%
  \renewcommand\makelabel[1]{\textbf{##1}\hfil}}}{\end{list}}
\makeatother

% ---- cover page --------------------------------------------------------------
\newcommand{\prelimcover}[3]{% #1 total marks, #2 duration, #3 date line
\thispagestyle{prelimcover}%
{\coverfont\fontsize{10}{12}\selectfont\linespread{1}\selectfont
\noindent\begin{minipage}[t]{\linewidth}
  {\bfseries \prelimtitle}\\[0.2ex]
  {\bfseries Compiled by Maximus Pung Jun}\\[1.6ex]
  Higher 2
\end{minipage}\par\vspace{2ex}
\noindent\rule{\linewidth}{1.2pt}\par\vspace{0.4ex}
\noindent{\Large\bfseries MATHEMATICS}\hfill{\Large\bfseries 9758}\par\vspace{1.2ex}
\noindent \prelimsubject\hfill\textbf{#3}\\
QUESTION PAPER\hfill\textbf{#2}\par\vspace{2.5ex}
\noindent\begin{tabular}{@{}p{46mm}l@{}}Additional Materials: & Writing paper\\ & List of Formulae and Results (MF27)\end{tabular}\par\vspace{2ex}
\noindent\rule{\linewidth}{0.6pt}\par\vspace{0.6ex}
{\bfseries READ THESE INSTRUCTIONS FIRST}\par\vspace{2ex}
\setlength{\parskip}{0pt}
Write your name and class on all the work you hand in.\\
Write in dark blue or black pen.\\
You may use an HB pencil for any diagrams or graphs.\\
Do not use staples, paper clips, glue or correction fluid.\par\vspace{2ex}
Answer \textbf{all} the questions.\\
Give non-exact numerical answers correct to 3 significant figures, or 1 decimal place in the case of angles in degrees, unless a different level of accuracy is specified in the question.\\
You are expected to use an approved graphing calculator.\\
Unsupported answers from a graphing calculator are allowed unless a question specifically states otherwise.\\
Where unsupported answers from a graphing calculator are \textbf{not} allowed in a question, you are required to present the mathematical steps using mathematical notations and not calculator commands.\\
You are reminded of the need for clear presentation in your answers.\par\vspace{2ex}
The number of marks is given in brackets [\ ] at the end of each question or part question.\\
The total number of marks for this paper is #1.
\vfill
\noindent\rule{\linewidth}{0.6pt}\par\vspace{0.4ex}
\centerline{This document consists of \textbf{\pageref{LastPage}} printed pages.}\par\vspace{1ex}
}\clearpage}

% ---- solutions header block (page 1 of the worked solutions; same cover style) ----
\newcommand{\prelimsolhead}[1]{% #1 total marks
{\coverfont\fontsize{10}{12}\selectfont\linespread{1}\selectfont
\noindent{\bfseries \prelimtitle}\\[0.2ex]{\bfseries Compiled by Maximus Pung Jun}\\[1.6ex]Higher 2\par\vspace{2ex}
\noindent\rule{\linewidth}{1.2pt}\par\vspace{0.4ex}
\noindent{\Large\bfseries MATHEMATICS}\hfill{\Large\bfseries 9758}\par\vspace{1.2ex}
\noindent \prelimsubject\hfill\textbf{2026}\\
WORKED SOLUTIONS\hfill\textbf{Total marks: #1}\par\vspace{1ex}
\noindent\rule{\linewidth}{0.6pt}\par}\vspace{4ex}}
